\documentclass[UTF8]{ctexart}
\usepackage{mathtools,wallpaper}
\usepackage{t1enc}
\usepackage{pagecolor}
\usepackage{enumitem}
\usepackage{hyperref}
\usepackage{graphicx}
\usepackage{listings}
\usepackage{amssymb}
\usepackage[dvipsnames]{xcolor}
\usepackage{float}
\usepackage{numerica}



\begin{document}

\section{L01-AI-Intro}
\subsection{基本概念介绍}
\begin{itemize}
    \item Multimodality
    
    多模态为模型引入了同时处理不同源，比如文字，音频，图像，视频等，的能力，以让模型获得更加丰富的理解。

    \item Agentic AI
    
    Agentic AI能够主动采取行动并在现实世界中发挥作用，不仅仅是响应查询。

    \item 人工智能
    
    人工智能是能够模拟人类智能某些方面的系统和机器。

    \item 传统编程
    
    人类专家明确地定义规则，程序一步步遵循这些规则

    \item 机器学习
    
    机器学习是人工智能的一个子集，它使系统能够通过算法从数据中学习，无需显示编程。人类无需编写规则，通过提供数据和相应标签让机器自动学习模式。

    \item 深度学习
    
    深度学习是机器学习的一个子集，深度学习使用多层神经网络去建模复杂的数据模式。

    \item 深度学习 vs 机器学习
    
    \begin{enumerate}
        \item 机器学习在数据量小的结构化数据训练效果好。深度学习需要大量非结构化的数据，比如图像音频文本。
        \item 机器学习使用简单的算法，比如线性回归，决策树，支撑向量机。深度学习使用复杂的卷积神经网络或者循环神经网络架构。
        \item 机器学习可以在标准的CPU训练，且耗时短。深度学习需要在高性能的GPU上训练，训练时间长。
        \item 机器学习模型更加透明且易解释。深度学习经常作为黑箱，解释困难。
    \end{enumerate}
    
    \item 自然语言处理
    
    自然语言处理研究计算机如何理解，生成和处理人类的语言。


\end{itemize}

\subsection{人工智能 ethical problems}
\begin{itemize}
    \item Job displacement and economic impact工作岗位替代和经济影响
    
    人工智能驱动的自动化可能会取代运输、客户服务等行业的工作人员。
    \item Bias and fairness
    
    人工智能系统可能会从训练数据中即成或放大偏见。
    \item Autonomous system and safety
    
    Who is responsible when AI causes harm?
    \item Misinformation and manipulation
    
    生成式AI能够创建出逼真的虚假信息。Ethical problem：虚假信息传播。
\end{itemize}

\section{L02-ML-workflow}
\subsection{机器学习类别}

机器学习包括有监督机器学习和无监督机器学习。有监督机器学习有回归问题和分类问题，分类问题中包括二分类问题和多分类问题。无监督机器学习包括聚类问题。有监督机器学习中，数据集包括其特征以及每个数据的标签；无监督机器学习数据集仅包含标签。


\begin{itemize}
    \item supervised learning监督学习
    
    监督学习通过构建基于一个或者多个输入来预测或估计输出的统计模型。训练集中包含特征和响应/标签，模型需要学习到通过新的数据预测响应的方法。

    \item classification分类问题
    
    预测离散类别/标签

    \item regression回归问题
    
    预测连续的数值

    \item unsupervised learning无监督学习
    
    无监督学习是指无需在带有标签的数据上训练模型就可以从数据上提取含义的统计方法。
    
    训练数据只有样本特征，没有标签。目的是找到数据中的模式、分组、结构或者内在关系。

    常见任务：降维，聚类，关联规则学习。

    \item clustering聚类
    
    在没有预定义标签的情况下对相似的数据进行分组
\end{itemize}

\subsection{机器学习工作流}

\subsubsection{确认问题类型}

\subsubsection{数据收集}

\begin{itemize}
    \item manual data collection手动数据收集
    
    surveys, interviews, manual labelling
    \item automated data collection自动化数据收集
    
    使用APIs从线上资源自动获取数据
    \item sensor or IoT data传感器或物联网数据
    
    从摄像头GPS或者温度传感器等物理设备获取数据
    \item Database and Log Data数据库和日志数据
    
    从已有的数据库或者系统日志提取数据
    \item public or open dataset
    
    公开数据集

    \item crowdsourcing众包

\end{itemize}

\subsubsection{Data Preprocessing数据预处理}
Data structuring, handle missing values, normalize or standardize features, and encode categorical variables.

\subsubsection{Feature Engineering特征工程}
Select, transform or create new features that can improve model performance.
\begin{itemize}
    \item 选择最相关的特征
    \item 将已有的特征通过取对数或者多项式或者interaction terms进行转换
    \item 从领域知识中创建新的特征
\end{itemize}

\subsubsection{模型选择}
根据问题和数据特征选择合适的模型
\begin{itemize}
    \item 分类：logistic regression, Naive Bayes, Nearest Neighbor, Discriminant analysis,decision trees, random forest, support vector machine, neural network
    \begin{enumerate}
        \item 逻辑回归：可以拟合一个预测属于一类或者另一类二分响应概率的模型
        \item 判别式分析：假设不同类别依据高斯分布(对于不同的类别，都有自己特别突出的特征)生成数据，通过寻找特征的线性组合对数据进行分类
        \item K近邻算法：根据数据集中待分类对象最邻近样本的类别来划分其类别
        \item 朴素贝叶斯：假设某类别中某个特征的存在与其他特征相互独立，基于数据属于某一类别的概率高低进行分类。
    \end{enumerate}
    \item 回归: Linear regression, Lasso regression, GLM, Gaussian Process, decision tree, random forests, gradient boosting
    \begin{enumerate}
        \item 线性回归：用于将连续响应变量描述(目标变量)为一个或者多个预测变量(特征)的线性函数
        \item 非线形回归
        \item 广义线性模型(generalized linear models)：使用将输入的线性组合拟合到输出的非线形函数上。
    \end{enumerate}
    \item 聚类: K-Means, hierarchical clustering层次聚类，DBSCAN密度聚类
    \item Data size数据量
    \begin{enumerate}
        \item 小的数据集，选择简单的模型如逻辑回归或者支撑向量机
        \item 大的数据集，选择复杂模型，如神经网络
        \item 如果小的数据集选择大的模型，那么会导致模型对训练集过拟合，因为小的数据集，特征较少，大模型的参数量很多，足以表示小的数据集中普通样本每一个特征的关系，也可以表示一些特殊样本负样本的特征的关系，导致过拟合。一个大的数据集选择简单的模型，那么模型会欠拟合，因为模型难以拟合数据特征的关系。
    \end{enumerate}
    \item Feature type特征类型
    \begin{enumerate}
        \item 数值数据:线性模型或者树模型
        \item 分类数据/混合数据:决策树，随机森林，梯度boosting模型
    \end{enumerate}
    \item 如果数据存在噪声或者缺失数据，则使用基于树的模型更加鲁棒
    \item 如果可解释性(intepretability)重要，使用线性回归、决策树或者逻辑回归
    \item 如果预测能力比可解释性重要，使用gradient boosting或者神经网络这些复杂的模型
    \item Evaluate Performance and Complexity
    \begin{enumerate}
        \item 从简单的模型开始训练
        \item 使用验证指标进行对比(RMSE，accuracy，F1)
        \item 若性能不足，使用复杂的模型
    \end{enumerate}
    \item 实验与tune调优
    \begin{enumerate}
        \item 使用cross-validation 比较模型
        \item hyperparameter tuning超参数调优(grid search，random search)
    \end{enumerate}
\end{itemize}
\subsubsection{Model Deployment模型部署}
将训练好的模型集成到生产环境中用于实际场景
\subsubsection{Monitoring and Maintenance监控与维护}
持续监控模型，在有新数据可用或者条件发生改变时对其进行更新

\section{机器学习基础}

\subsection{常见数据质量问题}
\subsubsection{Missing data数据缺失}
由于数据输入错误，设备故障或者未收集到的数值，一些特征是不完整的

\subsubsection{Inconsistent or unstandardized Data不一致或者不标准化的数据}
格式混杂，单位不一致，大小写不规范

\subsubsection{Imbalanced data distribution不平衡的数据分布}
一些模型占比过高，导致训练模型存在偏差
\subsubsection{Outliers and noisy data异常值和噪声数据}
干扰模型训练的极端值和错误值
\subsubsection{Redundant or Collinear Feature冗余或者共线特征}
重复和高度相关的特征会导致模型模型的不稳定性和过拟合。

\subsection{数据预处理}
\subsubsection{缺失值处理}
\begin{itemize}
    \item 填补或者移除缺失值
    \item Mean/median/KNN imputation(属于数字内插方法：使用简单的算法快速填补数据空缺)；interpolation插值方法(牛顿插值，拉格朗日插值，赫米特插值，三次样条插值)
\end{itemize}

\subsubsection{Normalization和standardization}
    移除尺度差异，加速模型收敛。
    \begin{enumerate}
        \item Normalization归一化(一般将数值缩放到0-1之间)
        归一化常常用于有数值范围的数据，比如像素值(0-255),一般用于图像数据处理、神经网络
        \begin{itemize}
            \item Min-Max Scaling最小最大缩放
            
            $x^{'}=\frac{x-x_{min}}{x_{max}-x_{min}}$
            
        \end{itemize}
        \item Standardize标准化(将数据标准化到均值为0，标准差为1)
        标准化一般用于位置边界的数据，用于线性模型、SVM、逻辑回归等
        \begin{itemize}
            \item Z-score Normalization
            
            $x^{'}=\frac{x-\mu}{\sigma}$
            $\mu$是特征的均值；$\sigma$是特征的标准差
        \end{itemize}
    \end{enumerate}

\subsubsection{Discretization离散化}
离散化将连续的数值变量通过分组他们的值到一个区间(bins)中来将连续的数值转换为离散的变量。

如果模型或算法需要离散的输入或者希望减少连续数据中的噪声或者简化连续数据中的关系。

\begin{itemize}
    \item Equal-Width Binning等宽分箱
    
    将数据范围划分为宽度相等的区间, 公式为$Bin width=\frac{x_{max}-x_{min}}{k}$，k为分箱的数量。区间除了最后一个是左闭右闭，前面都是左闭右开。
    \begin{enumerate}
        \item 优点
        
        实现简单
        \item 缺点
        
        如果数据分布不均，可鞥会导致空箱
    \end{enumerate}
    \item Clustering-based Binning聚类分箱
    
    使用聚类算法自动找到相似度高的箱。此方法是数据驱动的，能适应实际分布。
    \begin{enumerate}
        \item 优点
        
        自适应数据模式
        \item 缺点
        
        计算复杂性高
    \end{enumerate}
    \item Equal-frequency Binning等频分箱
    
    每个箱的样本数量相同，区间宽度随数据分布调整。按照数据排序后，均分样本到每个箱

    \begin{enumerate}
        \item 优点
        
        适配偏态数据(也就是非正态分布数据，数据有右侧拖尾，如薪资、房价)，对于偏态数据中位数通常比平均值更能代表数据的典型水平，避免极端值导致的分布不均
        \item 缺点
        
        每个箱的宽度不一样，边界无意义。
    \end{enumerate}

\end{itemize}

\subsubsection{Feature Encoding特征编码}
由于模型无法直接处理文本。特征编码将非数字的分类变量转换为数字形式的过程，以便机器学习算法能够对其进行处理。

\begin{itemize}
    \item label encoding标签编码
    
    每一个类别都被分配一个唯一的整数值。
    \begin{enumerate}
        \item 优点：
        
        简单，紧凑
        \item 缺点：
        
        整数值之间的关系暗示了顺序关系(2>1>0),对nominal categories(名义类别：是无内在顺序、无大小/等级关系的分类变量)产生误导。

        \item 适用对象：
        
        Tree model(决策树，XGboost)
    \end{enumerate}
    \item One-Hot encoding独热编码
    
    每个类别都表示一个二进制向量，只有一个元素为1，其余元素为0

    适用模型：线性模型，神经网络
\end{itemize}

\subsubsection{Data Balancing数据平衡}
在许多问题中，每个类别的样本数量是不均衡的。如果一个类别占主导地位，模型可能会忽略少数类别，从而在准确率上准确率较高，但对罕见事件招呼率很低。

\begin{itemize}
    \item oversampling过采样
    \begin{enumerate}
        \item 通过复制或者人工生成新样本来增加少数类样本的数量
        \item SMOTE(Synthetic minority Oversampling Technique):通过在少数类样本之间进行插值来合成少数类别样本
    \end{enumerate}
    \item Undersampling欠采样
    
    减少多数类样本的数量，使其与少数类样本数量匹配。

    简单且快速，但是可能会因丢失多数样本而丢失有价值的信息。
    \begin{enumerate}
        \item 优点
        
        能够保存所有数据集
        \item 缺点
        
        合成样本与真实数据模式有偏差，过度学习少数类细节，导致模型过拟合，泛化能力下降
    \end{enumerate}
    \item class-Weight Adjustment类别权重调整
    
    不改变数据的数量，提高样本数量少的类别的误分惩罚。

    公式：
    $\omega_{i}=\frac{N}{k*N_{i}}$

    i为类别，每一个i分配权重$\omega_{i}$。N是所有样本数量。$N_i$是第i类样本的数量。k是类别数量。

    \begin{enumerate}
        \item 优点
        
        能够保存所有数据集
        \item 缺点
        
        需合理调整权重参数否则可能引入新偏差
    \end{enumerate}
\end{itemize}

\subsubsection{Feature selection/extraction特征选择和特征提取}
减少维度，提升模型训练效率。移除不相关的冗余的特征。PCA主成分分析，LASSO岭回归，information gain信息增益

\begin{itemize}
    \item LASSO regression
    
    LASSO会将不重要的特征系数缩小到0，从而有效筛选关键特征。
    \item Information gain
    
    衡量每个特征对预测目标的贡献程度。保留信息增益最高的特征。
    \item Principal component analysis
    
    将现有特征转换或组合成新的特征。
\end{itemize}

\subsection{机器学习的3个基本组成部分}
机器学习的目标是使计算机能够从数据中自动学习模式，而不是依赖于明确的编程规则。

\subsubsection{Model模型}
模型表示输入x和输出y之间假设的数学关系。定义了x到y的映射函数$f(x;\theta)$,$\theta$是模型待学习的参数，不同的模型意味着不同的假设空间。

\begin{itemize}
    \item 对于机器学习任务，我们首先要定义输入空间x，和输出空间y，不同的机器学习任务之间的主要区别是输出空间y的形式。
    \begin{enumerate}
        \item 对于binary classification(二分类问题)：$y=\{0,1\}$或者$y=\{-1,1\}$
        \item 对于多分类问题(N分类问题)：$y=\{1,2,...,N\}$
        \item 回归问题:$y=\mathbb{R} $
        \item 输入空间x通常表示样本的特征空间
        \item 可以表示为$p_r(y|x)$或者$y=g(x)$
    \end{enumerate}
    \item Linear Model线性模型
    
    线性模型假设输入x和输出y是线性关系的，其假设空间为参数化的线性函数族。

    对于特征向量$x=[x_1,x_2,...,x_d]^T$和权重向量$w=[w_1,w_2,...,w_d]^T$，b是偏置,那么线性模型预测值可以写为如下公式：

    $\hat{y}=w^Tx+b$

    \item Nonlinear Models非线形模型
    
    $\hat{y}=f(x;\theta)$

\end{itemize}

\subsubsection{Learning Criterion学习准则/Objective Function目标函数/loss Function损失函数}
目标函数定义了模型性能的好坏，衡量预测输出与真实输出之间的差异(discrepancy)。包括均方误差，交叉熵损失等。

\begin{itemize}
    \item Binary classification二元分类 
    
    二元分类问题，每个样本可以分为正类和负类。因此预测结果可以分成4种(第一位为模型是否预测正确，第二位为模型预测的结果)：
    \begin{enumerate}
        \item True positive(真阳性)：实际为正例，模型正确预测其为正例的结果。
        \item False negative(假阴性)：实际为正例，模型错误预测其为负例的结果。
        \item False positive(假阳性)：实际为负例，模型错误预测其为正例的结果。
        \item True negative(真阴性)：实际为负例，模型正确预测其为负例的结果。
    \end{enumerate}
    \item confusion matrix混淆矩阵
    \begin{enumerate}
        \item Accuracy(准确率) 
        准确率衡量所有样本中，被正确分类的样本所占的比例。$Accuracy=\frac{TP+TN}{TP+TN+FP+FN}$

        但在数据不平衡时，效果较差。

        \item Precision(精确率)
        
        精确率衡量的是预测为阳性的样本中，实际为阳性的比例。$Precision=\frac{TP}{TP+FP}$

        高精度意味着当模型预测为阳性时，它通常是正确的。精确率侧重于减少假阳性。
        \item Recall(召回率)
        
        召回率衡量的是模型正确识别出实际正样本所占的比例。$Recall=\frac{TP}{TP+FN}$

        高召回率意味着模型成功捕捉到了大多数的阳性样本，召回率为了减少假阴性。

        \item F1-score(F1分数)
        
        F1分数代表精确率和召回率达调和平均值，提供了一个平衡这两个指标的单一分数。$F1-score=2*\frac{Precision*Recall}{Precision+Recall}$
    \end{enumerate}

    \item K-Fold Cross-Validation(K折交叉验证)
    \begin{enumerate}
        \item K折交叉验证将数据集随机划分为k个大小相等的子集；在每一轮验证中选择一个子集作为测试集，其余k-1个子集作为训练集；此过程重复k次，确保每个子集都作为测试集使用1次；最后模型性能通过对所有k轮的结果取平均值评估。
        \item 一般k=10
    \end{enumerate}

    \item 分类算法中的Training error(训练误差) and Generalization error(泛化误差)
    \begin{enumerate}
        \item 训练误差：模型对现有的数据集的拟合程度，衡量模型在已见过的数据上的表现。
        \item 泛化误差：模型对新的未见过的数据进行正确分类的能力，衡量模型在训练集之外的泛化能力。
    \end{enumerate}

    \item Underfitting欠拟合 and Overfitting过拟合
    \begin{enumerate}
        \item 欠拟合：模型无法捕捉训练集数据的潜在模式。可以通过增加特征数量，选择更相关的特征以及使用更复杂的模型结果等方法来解决。
        \item 过拟合：模型在训练数据上过于特化，在新的未见过的样本上表现不佳，导致训练误差持续下降，泛化误差开始上升。
    \end{enumerate}

\end{itemize}

\subsubsection{Optimization algorithm优化算法}
优化算法通过调整模型参数来最小化/最大化学习准则。包括：梯度下降，随机梯度下降，动量法，Adam，AdamW


\section{Artificial Neural network(ANN)人工神经网络}

\subsection{从Neurons(神经元)到Perceptrons(感知机)}

感知机是线性二元分类器，常用于监督学习，要求样本是线性可分的。公式为:$y=f(\sum\limits_{i=0}^{n} w_{i}x_{i}+b)$

x是输入向量，w是权重向量(突触)，b是bias(偏置)，f函数是activation function激活函数，Y是输出信号。

\begin{itemize}
    \item 输入值
    \begin{enumerate}
        \item 输入值是每个样本的features，样本在训练过程中按batch提供
        \item 更大规模更多样化的输入特征通常会提升模型性能。
    \end{enumerate}
    \item Weights and Bias
    \begin{enumerate}
        \item 权重和偏置是可学习的参数。网络通在前向传播和反向传播的训练过程中，这两个参数会被调整以达到期望的值或输出。
        \item 权重用于确定每个特征的重要性
        \item 偏置使模型能够移动激活函数。
    \end{enumerate}
    \item weighted sum加权和
    \begin{enumerate}
        \item 将样本的各个特征与各自权重相乘，$z=\sum_{w_j}^{x_j}+b$,w是权重，x是输入特征
        \item 它是激活函数的输入
        \item 在训练过程中通过反向传播不断更新，来最小化损失。
    \end{enumerate}
    \item 激活函数activaton function
    \begin{enumerate}
        \item 使用step function阶跃函数：如果加权和结果大于0，就为1；如果小于0，就为0.
    \end{enumerate}
    \item 前向传播
    
    输入样本特征值，乘以权重，加权和，输入激活函数，得到输出信号

    \item 反向传播
    
    损失函数对参数w和b进行求偏导以最小或最大化损失，使用偏导结果和学习率更新参数

    \item 感知器的局限性
    \begin{enumerate}
        \item Single-layer structure单层架构
        感知器是单层架构，只有一层权重。它完全依赖所提供的输入特征，没有能力发现或学习相关的中间表示。
        \item No ability to handle Non-linearity没有能力处理非线性问题
        感知器使用线性激活函数，它无法在学习过程中引入非线形。
    \end{enumerate}

    \item Multi-Layer Perceptron (MLP多层感知机)
    
    多层感知机具有多层结构，非线形激活函数和各种损失函数，能够处理复杂的非线形问题和多分类问题。
\end{itemize}

\subsection{Activation function激活函数}
激活函数引入了非线性，使模型能够学习超越线性关系的复杂模式。

\begin{itemize}
    \item Sigmoid函数
    \begin{enumerate}
        \item sigmoid函数产生0到1之间的输出，适用于probabilistic interpretations.
        \item $\sigma(x)=\frac{1}{1+e^{-x}}$
    \end{enumerate}
    \item tanh函数
    \begin{enumerate}
        \item tanh函数产生-1到1之间的输出，常被使用于中心化数据(也就是数据均值为0)。
        \item $tanh(x)=\frac{e^{x}-e^{-x}}{e^{x}+e^{-x}}$
    \end{enumerate}
    \item ReLU(Rectified Linear Unit整流线形单元)
    \begin{enumerate}
        \item 对负输入输出零，对正输入输出线性值。
        \item $\operatorname{ReLU}(x) = \begin{cases} x, & \text{if } x > 0 \\ 0, & \text{otherwise} \end{cases}$
        \item $\text{ReLU}(x) = \max(0, x)$
    \end{enumerate}
    \item Leaky ReLU
    \begin{enumerate}
        \item $\operatorname{ReLU}(x) = \begin{cases} x, & \text{if } x > 0 \\ 0.1x, & \text{otherwise} \end{cases}$
        \item $\text{ReLU}(x) = \max(0.1x, x)$
    \end{enumerate}
    \item ELU
    \begin{enumerate}
        \item $\operatorname{ELU}(x) = \begin{cases} x, & \text{if } x \geqslant  0 \\ \alpha(e^{x}-1), & x<0 \end{cases}$
    \end{enumerate}
    \item Softmax
    \begin{enumerate}
        \item $Softmax(x)=\frac{e^{x}}{\sum\limits_{i}e^{x_{i}}}$
    \end{enumerate}
    \item GELU
    \begin{enumerate}
        \item $0.5x(1+tanh(\sqrt{\frac{2}{\pi}}(x+0.044715x^3)))$
    \end{enumerate}
    \item Maxout
    \begin{enumerate}
        \item $max(\omega^T_1x+b_1,\omega^T_2x+b_2)$
    \end{enumerate}
    \item 选择的激活函数与输出范围和目标值类型或该层的作用相关
    \begin{enumerate}
        \item 二元分类：Sigmoid函数
        \item 多分类问题：Softmax函数，值域为0到1，结果和为1
        \item 隐藏层：整数或者混合实数值，选择ReLU、Leaky ReLU，GELU
        \item Recurrent neural network: 状态取值(-1,1),使用Tanh
    \end{enumerate}
\end{itemize}

\subsection{Loss function损失函数}
损失函数是非负实值函数，用于量化模型预测与真实标签之间的差异。

\begin{itemize}
    \item 0-1 loss function
    
    0-1损失函数衡量预测是正确的还是错误的.$\operatorname{L}((y,\hat{y})) = \begin{cases} 0, & \text{if } \hat{y} = y \\ 1, & \hat{y}\neq y \end{cases}$

    y是真实标签，$\hat{y}$是预测标签。
    
    \begin{enumerate}
        \item 缺点
        \begin{itemize}
            \item Non-differentiable:损失不连续且导数为0，无法直接通过基于梯度的方法直接优化。
            \item 无法直接反应模型对预测结果的置信度的高低。
        \end{itemize}
    \end{enumerate}

    \item Cross-Entropy Loss Function交叉熵损失函数
    
    交叉熵损失衡量两个概率分布之间的差异，量化了预测概率分布与真实分布的匹配程度。
    y是(一般通过独热向量编码)真实标签，$\hat{y}$是(通过sigmoid或者softmax)预测概率
    \begin{enumerate}
        \item binary cross-entropy二元交叉熵/Logistic损失函数
        
        $L=-\frac{1}{N}\sum\limits_{i=1}^{N}[y_{i}log(\hat{y_{i}})+(1-y_{i})log(1-\hat{y_{i}})]$
        \begin{enumerate}
            \item $y_{i}是真实标签，y_{i}\in \{0,1\}$
            \item $\hat{y_{i}}是预测概率，\hat{y_{i}}\in (0,1)$
            \item 在sigmoid之后使用
            \item Log损失是0-1损失的平滑近似，可以产生概率分布。
        \end{enumerate}
        \item categorical cross-entropy分类交叉熵
        
        $L=-\frac{1}{N}\sum\limits_{i=1}^{N}\sum\limits_{c=1}^{C}y_{i,c}log(\hat{y_{i,c}})$
        \begin{enumerate}
            \item i是样本的索引，c是类别
            \item $y_{i,c}$，如果c是正确的类别，则为1；否则为0
            \item $\hat{y_{i,c}}$，表示第i个样本属于第c类的概率
        \end{enumerate}
    \end{enumerate}
    \item Kullback-Leibler Divergence(KL散度)
    
    $Loss=\sum\limits_{i}P(x_{i})log(\frac{P(x_{i})}{Q(x_{i})})$

    \item Hinge Loss(铰链损失函数)
    
    $Loss=max(0,1-\hat{y}y)$
    \begin{enumerate}
        \item 适用于支撑向量机模型
        \item 不同类别中存在间隔，关注误分样本和边界样本
    \end{enumerate}

    \item Mean Square Error均方误差
    
    $MSE=\frac{1}{n}\sum\limits_{i=1}^{n}(y_{i}-\hat{y_{i}})^{2}$
    \begin{enumerate}
        \item 适用于回归任务
        \item 计算目标值与预测值的平方平均值
        \item 对异常值敏感
    \end{enumerate}

    \item Mean Absolute Error平均绝对误差
    
    $MAE=\frac{1}{n}\sum\limits_{i=1}^{n}|\hat{y_{i}}-y_{i}|$
    \begin{enumerate}
        \item 采用目标值与预测值的绝对差值，相比MSE对异常值鲁棒性更强。
    \end{enumerate}

    \item Root Mean Square Error均方根误差
    
    $MSE=\sqrt{\frac{1}{n}\sum\limits_{i=1}^{n}(y_{i}-\hat{y_{i}})^{2}}$
    \begin{enumerate}
        \item 适用于回归任务
    \end{enumerate} 

    \item Huber Loss
    
    $L_{\delta}(y, f(x)) = \begin{cases} \frac{1}{2}(y - f(x))^2, & |y - f(x)| \leq \delta \\ \delta|y - f(x)| - \frac{1}{2}\delta^2, & |y - f(x)| > \delta \end{cases}$
    
    \begin{enumerate}
        \item 融合MAE和MSE的优势，小误差时使用二次函数损失，大误差使用线性函数损失
    \end{enumerate}
\end{itemize}

\section{Deep Neural Networks(DNN)深度神经网络}
\subsection{深度神经网络}
\begin{itemize}
    \item 训练过程
    \begin{enumerate}
        \item 前向传播：计算模型输出和损失
        \item 反向传播：计算计算模型梯度，使用链式法则更新模型参数。
        
        \begin{enumerate}
            \item $\bigtriangleup w_{ji}=-\eta \frac{dE}{d w_{ji}}E$\end{enumerate}
            \item 神经网络有很多参数，因此损失函数是一个由这些参数组成的多变量函数，我们要计算每个参数的偏导，以了解每个参数如何影响损失
            \item 如果想知道每一层的参数如何影响损失，我们需要从最后一层输出层一步一步向前求每个参数的偏导数(占用大量计算资源)
        \end{enumerate}
        
        
        \item 参数更新：使用梯度和学习率自适应权重
        \item 循环以上过程
    
    \item universal function approximation Theorem通用函数逼近定理
    
    给定任意连续函数f(x),如果一个2层的神经网络有足够的隐藏单元，那么存在一种权重选择，使其能够以任何期望的精度紧密逼近任何连续函数f(x)

    \begin{enumerate}
        \item 大量神经元会导致过拟合风险，需要一些早停或者正则化策略
        \item Automatic differentiation自动微分
    \end{enumerate}

    \item Scaling laws缩放定律
    \begin{enumerate}
        \item Model size：更大的模型/更多的模型参数通常在任务上表现更好
        \item Data Volume：使用更多的数据训练，通常会提高准确性
        \item Computation计算能力：增加计算资源可以训练更深更大的网络
    \end{enumerate}

    \item 深度神经网络
    
    $z_{i}^{k}=g(\sum\limits_{j}W^{(k-1,k)_{i,j}}z^{(k-1)}_{j}+b^{(k)}_i)$
    \begin{enumerate}
        \item 第k层的第i个神经元等于连接第k-1层第j个神经元与第k层第i个神经元的权重乘以第k-1层的神经元输出之和加上偏置的数值，最后输入激活函数g
    \end{enumerate}
\end{itemize}

\subsection{Deep Learning深度学习}
\subsubsection{Data preparation数据准备}‘
数据被分为训练集，测试集和验证集。训练集用于模型训练。验证集一般用于超参数选择，测试集用于模型泛化能力测试。
\begin{itemize}
    \item Data Collection \& Quality
    \begin{enumerate}
        \item 确保你有足够多与你要解决问题相关的较高质量的数据
    \end{enumerate}
    \item Data preprocessing
    \begin{enumerate}
        \item 归一化/标准化：对数据进行标准化，确保模型更快更稳定地收敛
        \item Data Augmentation数据增强：对于图像任务，可以通过transformation(比如：旋转，翻转，缩放)创建更多样的训练数据
    \end{enumerate}
    \item Handling Imbalanced Datasets
    \begin{enumerate}
        \item 考虑对少数类进行过采样，多数类进行欠采样，
        \item 使用带权重的交叉熵损失，focal loss的相关的损失函数
    \end{enumerate}
\end{itemize}

\subsubsection{Model Design}
\begin{itemize}
    \item Architecture架构
    \begin{enumerate}
        \item 基于问题选择什么类型的模型(CNN适合图像数据训练；RNN和Transformer适合时序数据)
        \item 选择层数以及层类型(convolutional卷积层, fully connection全连接层, recurrent循环层)
        \item 越深的模型可以捕捉到越复杂的模式，但是很难训练
    \end{enumerate}
    \item Layer Sizes层大小
    
    层大小表示the number of neurons in each layer. 越多的神经元会提高模型能力，但对与小数据集会导致过拟合的风险。
    \begin{enumerate}
        \item Funnel Rule漏斗规则
        
        输入->large hidden size->smaller hidden size->output

        \item 数据大小 VS 模型大小
        
        \begin{itemize}
            \item 参数数量<10*样本数量
            \item 使用regularization(dropout或者L2)正则化
            \item 从简单的网络开始，向大网络扩展
        \end{itemize}
    \end{enumerate}

    \item 激活函数
    
    影响模型学习非线形关系的方式。

    \item 输出层和损失函数
    
    \begin{enumerate}
        \item 根据任务选择输出层size。分类任务：softmax；回归：线形层
        \item 损失函数：损失函数取决于具体任务。分类问题：交叉熵损失；回归：使用均方误差
    \end{enumerate}

    \item 机器学习中的Convergence
    
    收敛指训练过程中模型参数和损失函数值趋于稳定状态的过程。

    \begin{enumerate}
        \item 损失值随着迭代不断减小，参数更新变小。也就是模型学到了足够多的知识。
    \end{enumerate}

    \item 机器学习中的divergence
    发散指训练过程没有朝着稳定改进的方向发展，而是训练过程变得不稳定
    \begin{enumerate}
        \item 训练损失没有下降反而上升了
        \item 损失大幅度波动
        \item 参数更新变化很大，且不稳定(overshooting超调)
    \end{enumerate}

\subsubsection{Training Process训练过程}

\begin{itemize}
    \item Optimizer优化器
    
    优化器通过调整模型的参数来最小化损失函数。比如SGD随机梯度下降，RMSprop，Adam等。

    \begin{enumerate}
        \item 训练的目标是通过优化器更新模型参数，最小化损失函数L.$\theta\leftarrow \theta-\eta \frac{dL}{d \theta}$
        \item $\eta$是learning rate学习率 
        \item $\frac{dL}{d \theta} $是损失函数对于参数的偏导数
        \item 优化器类型
        \begin{itemize}
            \item Stochastic gradient descent(SGD随机梯度下降)
            
            \begin{enumerate}
                \item 根据损失函数对参数的梯度更新权重
                \item $w_{nex}=w_{old}-\eta \bigtriangledown L(w)$
                \item $\eta$是学习率；$\bigtriangledown L(w)$是梯度
            \end{enumerate}

            \item Momentum动量法
            
            \begin{enumerate}
                \item 使用上一步的梯度平滑权重更新并且可以预防收敛到局部最小值。
                \item 为什么叫动量法，我们使用之前的速度的动量来冲过当下的局部极小值
                \item $v_{t}=\beta v_{t-1}+(1-\beta)\bigtriangledown L(w)$,$w_{new}=w_{old}-\eta v_{t}$
                \item 如果$\beta$小，那么当前的偏导对当前的梯度影响大；若$\beta$大，那么上一步梯度对当前的梯度影响大
            \end{enumerate}

            \item RMSprop(root mean square propagation均方根传播)
            
            \begin{enumerate}
                \item $v_{t}=\beta v_{t-1}+(1-\beta) (\bigtriangledown L(w))^2$；
                \item $w_{new}=w_{old}-\frac{\eta}{\sqrt{v_{t}+\epsilon }}\bigtriangledown L(w)$
                \item 将学习率除以梯度的移动平均值。通过除以每一个参数的梯度，控制每一个参数的步长/学习率，如果一个参数上一步的更新梯度很大的话，容易导致更新不稳定，那么学习率除以$v_t$会导致学习率变小，避免更新补偿太大导致震荡。
                \item 若梯度小，参数更新较慢，那么学习率除以梯度导致学习率变大，加速收敛。
                \item 在计算梯度是加平方：消除方向的影响，关注参数更新的波动程度；在求梯度时减少计算量。
                \item 可以处理梯度消失的问题
            \end{enumerate}

            \item  Adaptive Moment Estimation(Adam，自适应动量估计)
            \begin{enumerate}
                \item $m_{t}=\beta m_{t-1}+(1-\beta)\bigtriangledown L(w)$
                \item $v_{t}=\beta v_{t-1}+(1-\beta)(\bigtriangledown L(w))^2$
                \item Weight Update:$w_{new}=w_{old}-\eta\frac{m_{t}}{\sqrt{v_{t}+\epsilon}}$
                \item 有RMSprop的优点和Momentum的优点
            \end{enumerate}

        \end{itemize}
    \end{enumerate}
    \item Learning rate学习率
    
    学习率控制优化器在梯度方向上的步长大小
    \begin{enumerate}
        \item 太高的学习率可能会阻碍收敛，过低的学习率导致训练缓慢
        \item 可以使用学习率调度(步进衰减，指数衰减，余弦退火)或者自适应学习率
    \end{enumerate}

    \item Batch size批大小
    
    一个batch size表示每一次梯度更新训练的样本数量

    \begin{enumerate}
        \item 较小的batch size可以加快泛化速度，更好地适配新数据。因为小批量样本随机性强，导致每次更新参数时可能学习到更多新知识，帮助模型跳出局部最优，避免过度依赖训练数据细节。
        \item 大batch size加快计算速度，且梯度更稳定。大批量中，GPU一次性处理多个样本，充分利用其并行计算能力；大量样本梯度平均后波动更小，不会频繁震荡。但是大样本会覆盖大多数的训练集，导致过拟合。
    \end{enumerate}

    \item Number of Epochs训练轮次
    
    一个训练轮次是使用训练数据集中所有样本进行一次训练的完整训练周期，可以包含多次batch size。

    \begin{enumerate}
        \item 如果epoch数量少，模型可能不能学习到数据中的潜在模式，导致欠拟合
        \item 若epoch过大，模型可能会深究训练集的细节，导致过拟合，在未见过的新的数据集上泛化能力不佳。
        \item 理想的epoch次数通过记录模型在每个epoch训练完之后在验证集上的性能来确定
    \end{enumerate}

    \item Regularization Techniques正则化技术
    
    正则化是用于减少神经网络过拟合的技术，dropout，L1正则，L2正则，批归一化。

    \begin{enumerate}
        \item Dropout:在训练过程中，每次前向传播时每个神经元按照一定概率(20\%-50\%)随机丢弃。减少有效神经元数量，同时反向传播时这些神经元值的梯度值也为0，不影响参数更新，可以防止过拟合。
        \item L2正则(Weight Decay权重衰减):会在损失函数中添加一个惩罚项，抑制较大的权重值。$L_{total}=L_{original}+\lambda\sum\limits_{i}w_{i}^{2}$。$\lambda$是正则化因子。减小模型权重，让模型更平滑，在加权和中，权重减小，即使x有变化，输出y的波动也比较小。
    \end{enumerate}

    \item evaluation metrics评估指标
    
    选择合适的评估指标，accuracy,precision, recall, F1 score, AUC-ROC.

    \begin{enumerate}
        \item ROC Curve(Receiver Operating Characteristic Curve)
        
        \begin{enumerate}
            \item Y轴：真阳性率True Positive Rate=TPR=Recall=$\frac{TP}{TP+FN}$
            \item X轴：假阳性率False Positive Rate=FPR=$\frac{FP}{TP+TN}$
            \item ROC曲线上的每一个点对应不同分类阈值的(假阳性率,真阳性率)。对于每一个阈值，求其真阳性率和假阳性率，若曲线在对角线之上，则模型效果可以。
        \end{enumerate}

        \item AUC(Area Under tne Curve)
        
        AUC量化了模型区分正类和负类的整体能力，曲线下面积越接近1，模型效果越好，越接近0.5模型效果越差
    \end{enumerate}

\end{itemize}

\end{itemize}

\section{Convolutional neural networks(CNN卷积神经网络)}

\subsection{神经网络特征学习}
\subsubsection{Hierarchical feature learning分层特征学习}
\begin{enumerate}
    \item 每一层都学习不同层次的特征
    
    \begin{itemize}
        \item Shallow layers浅层(靠近输入)学习low-level patterns低级模式。（边缘，角点，颜色梯度）
        \item Middle layers将简单的模式组合成更有意义的结构。（纹理和形状）
        \item deep layer学习更高级的抽象特征。（物体部件）
    \end{itemize}

    \item 反向传播使网络能自动学习特征
    \begin{itemize}
        \item 网络做出预测
        \item 将预测结果与真实标签进行比较，得到损失
        \item 损失被反向传播，告诉每一层调整权重，使输出更接近真实情况
        \item 每一层都会更新其参数，以检测能够减少损失的模式
    \end{itemize}
    
    \item 传统全连接层在处理图像时面临的两个问题
    \begin{enumerate}
        \item 参数过多：对于100*100*3的图像，第一个隐藏层的每一个神经元都需要30000个独立连接。导致数据量大，且容易过拟合
        \item lack of local invariance缺乏局部不变性：图像对象具有局部特征，这些特征在缩放旋转移位等操作下基本不变。若不进行大量数据增强，难以捕捉到这种不变性。
    \end{enumerate}
\end{enumerate}

\subsection{CNN}
卷积神经网络是深度前馈神经网络，具有local connections局部连接，weight sharing权值共享，spatial invariance空间不变性。

\begin{itemize}
    \item 深度学习中卷积运算
    \begin{enumerate}
        \item 深度学习中卷积运算(torch库)直接进行滑动点积，不进行翻转(numpy库)。因为深度学习中参数是可学习的，不是固定的滤波器算子；省去翻转操作可以减少一次不必要的计算步骤。
    \end{enumerate}
    \item Padding填充
    \begin{enumerate}
        \item border effect边界效应：图像的中心像素会参与大量的卷积运算，图像的边缘位置参与卷积运算次数较少。比如中心位置的像素能被5*5的卷积在不同的卷积元素位置上覆盖，但是边缘像素只会覆盖一次。
        \item 造成图像信息损失：边缘像素的特征无法充分参与卷积计算，导致输出特征图中边缘信息丢失
        \item 图像尺寸变小：卷积后的特征图尺寸可能比输入小
        \item 0-padding：在输入图像的外围补0，减轻边界效应，同时增加生成的特征图的空间维度
    \end{enumerate}
    \item Pooling Layer池化层
    
    卷积层之后是池化层。
    \begin{enumerate}
        \item 池化目标
        \begin{itemize}
            \item Spatial dimension reduction(减小空间维度)
            
            池化从过滤区域提取一个值来表示这个区域，因此会减小矩阵的大小，并降低网络计算复杂度。
            \item Feature generalization(特征泛化)
            
            输入的微小变化和偏移将不会影响池化后的特征图，这是的网络对微小变化具有不变性。
            \item Increased receptive field(增大感受野)
            
            感受野是CNN中某一层神经元能看到的原始输入图像的区域，若卷积核为3*3，那么感受野为图像中的3*3区域。再使用2*2的池化层，就会变成6*6的感受野。

            通过对特定邻域像素窗口进行池化，网络能从更大区域捕捉模式，增大感受野。
        \end{itemize}
        \item MaxPooling最大池化
        
        最大池化选择每个池化窗口的最大值并且保留窗口中的主要特征。
        \begin{itemize}
            \item 帮助卷积神经网络关注最重要的模式
            \item 让模型对噪声具有鲁棒性
        \end{itemize}
        \item AveragePooling平均池化
        
        平均池化选择每个池化窗口的平均值并且保持整体信息的平滑。
        \begin{itemize}
            \item 保留更多的背景信息
            \item 当提取图像中像素强度分布时有用
        \end{itemize}

        \item 池化缺点
        \begin{enumerate}
            \item 池化会损失一些信息，尤其是一些较弱但是重要的特征(边缘/角点)
            \item 池化是一种不可学习的操作。
        \end{enumerate}
    \end{enumerate}
    \item stride步长(strided convolution)
    
    步长S>1会跳过输入中的某些位置，降低输出特征图的空间分辨率(下采样)。

    优点：
    \begin{enumerate}
        \item 由于卷积是可学习的，虽然步长仍然是下采样，但是学习到了如何以最优方式进行下采样
    \end{enumerate}

\subsection{训练CNN}
    \item 构建CNN网络
    \begin{enumerate}
        \item 输入，卷积，激活函数，最大池化，卷积，激活函数，最大池化，全连接层(线性层，激活函数，线性层)，全连接层，输出
    \end{enumerate}
\end{itemize}

\section{natural language processing(NLP自然语言处理)介绍}

\subsection{NLP概述}
\begin{itemize}
    \item NLP的定义：理解和生成人类语言的计算方法
    \item 挑战：将语言的semantics(语义)和pragmatics(语用)融入计算机理解非常困难。
    \begin{enumerate}
        \item Ambiguity(语言歧义，相同语言不同语义)
        \item Paraphrasing(释义相同语义的不同表达方式)
        \item Socio-Linguistic Variation(社会语言变体：同一语言因使用场景、人群不通，在表达上存在差异)
        \item Multi-Language and Code-Switching(多语言混合和缩写)
    \end{enumerate}
\end{itemize}

\subsection{NLP Upstream tasks上游任务}

上游任务是指为下游应用任务打基础的核心预处理任务，他们不直接解决实际业务问题，而是通过处理原始文本、提取基础语言特征，为后续复杂应用提供支撑。

\subsubsection{Preprocessing Tasks:Tokenization分词 \& Sentence spliting分句}

\begin{itemize}
    \item Tokenization分词:将文本拆分为称为token的小单元。
    \item N-grams:由n个连续单词组成的序列，用于捕捉局部语境并预测语言模式
    \item 句子拆分：将文本分割成句子
    \begin{enumerate}
        \item 拆分的句子在句法上是完整的，并且代表完整的思想
        \item 拆分的句子可以支持需要小文本单元的任务，比如机器翻译
        \item 拆分的句子支持并行处理的任务。句子分割后，多个句子是相互独立的，可同时分配个多个计算核心进行处理。
    \end{enumerate}
\end{itemize}

\subsubsection{Language Modeling(LM语言建模)}

\begin{itemize}
    \item 定义：基于前文已有的单词，预测序列中的下一个单词，以捕捉语言的架构和模式，理解语境。
    \item 一个句子的概率，可以被定义为给定先前符号每一个符号概率的乘积。
    \item 语言模型的类型：
    \begin{enumerate}
        \item Statistical Models统计模型：概率预测
        \item Neural Models神经模型：深度学习模型(RNN，Transformer)
    \end{enumerate}
\end{itemize}

\subsubsection{Part of Speech(POS) Tagging词性标注}
\begin{itemize}
    \item 目标：由于词语存在歧义，因此要确定每个词的syntactic categories(语法类别:noun,verb...)，以及additional properties(附加属性：tense时态)
    \item Annotated Corpus标注语料库：Pre-labeled data for training(用于训练的预标注数据)
    \item 词句的特征:
    \begin{enumerate}
        \item 基本特征：suffix/prefix前后缀，sentence position句子位置，capitalization大写，punctuation标点 来确定当前词的词性
        \item context window(上下文窗口)：融合相邻token的特征来决定当前词的词性
        \item Previous POS Tags(先前词性标签)：将先前词语的词性预测结果当作特征辅助当前单词的词性判断
    \end{enumerate}
    \item 机器学习预测POS Tagging
    \begin{enumerate}
        \item 算法：逻辑回归，决策树，随机森林，支撑向量机
        
        将当前token，前一个token，后一个token的的特征映射到词性标签，每个词元单独分类。
        \begin{enumerate}
            \item 先对原始文本分局，再进行分词。标签编码：将词性标签变成整数索引，便于模型计算损失，构建成标签映射字典，训练后用于反向解码。模型输入单token特征向量(基础特征+语境窗口特征：当前词的特征向量=自身词形特征+前词的词形特征+后词词形特征)。经过模型，预测当个模型的词性标签概率分布，求损失。
            \item 模型通过学习标注好的数据集中数据规律(比如：I会标注成PRP表示人称代词，like标注成VBP表示动词原形等)，建立特征组合与词性标签的对应关系。预测时，只依赖自身及前后词的特征，不考虑其他词的词性标签结果。
        \end{enumerate}
        \item word clustering单词聚类：将相似的单词分成k组以减小特征规模
        \item sequence labeling序列标注：
        \begin{enumerate}
            \item Conditional Random Fields(CRF条件随机场)基于大概率出现的标签模式标注序列。CRF参考词性标签应有的合理排列规律来输出争端序列的标注结果。
            \item 找到最大标签序列y$argmax_{y}p(y|x)$使整段序列一次性完成标注
        \end{enumerate}
    \end{enumerate}
    \item Named Entity Recognition(NER命名实体识别)
    
    识别文本中的实体，并将其分类到预定义的类别中。

    识别非结构化的文本中的关键信息来提供上下文和结构。


    \item Coreference Resolution共指消解
    
    识别文本中指向同一实体的不同表达。

    \item Word Sense Disambiguation词义消歧
    
    词义消歧指根据上下文语境，为文本中的多义词确定其在该场景下的唯一真实含义。

    \item Semantic Role labeling
    
    语义角色标注/论元识别与标注指为句子中的核心谓语/动词等找出去关联的各类语义角色，并标注每个角色与谓词的语义关系。
\end{itemize}

\subsection{NLP 应用}

\begin{itemize}
    \item Text classification(TC)
    
    文本分类任务自动将文本分类成预定义的标签或主题。

    \item Machine Translation(MT)
    
    机器翻译任务自动将文本从一门语言转为另一门语言。

    方法
    \begin{enumerate}
        \item Rule-based:依赖于linguistic rule(语言规则)和grammar(语法)
        \item Statistical:使用bilingual text corpora(双语语料库)进行概率预测。统计性概率翻译让模型从双语语料库中统计规律，统计语料中中文词或短语与英文词或短语的共现概率。统计语料中中文句子结构对应英文句子结构的转换概率。
        \item Neural:使用深度学习模型
    \end{enumerate}

    \item Information Retrieval(IR)
    
    信息检索是基于用户的询问，在大数据库中发现相关信息的过程。

    过程
    \begin{enumerate}
        \item Indexing索引：组织数据以快速访问
        \item Query processing询问处理：解析用户输入以检索相关结果
        \item Ranking：按照查询的相关程度排序结果
    \end{enumerate}

    方法
    \begin{enumerate}
        \item Boolean Retrieval：使用布尔运算符进行精确匹配
        \item probabilistic Model:使用概率对文档进行排序 (BM2.5/TF-IDF)。不理解文本语义，仅从词频、文档长度、查询词重要性等表面特征计算查询词出现在相关文档中的概率，以此作为相关性评分。
        \item 神经网络：采用深度学习根据相关性对结果进行排序。
    \end{enumerate}



    \item Question Answer
    
    问答系统用于自动回答自然语言提出的问题。

    \begin{enumerate}
        \item closed-domain(封闭领域)：在一个特定主题领域内回答问题
        \item Open-domain(开放领域)：可以回答广泛主题的问题
        \item Retrieval-based基于检索(extractive抽取式)：查找相关的文档或段落并且抽取答案
        \item Generative Model生成式模型(abstractive)：使用神经网络模型直接从知识中生成答案。(神经网络通过大规模的训练参数，就是学习到的内化知识)
    \end{enumerate}

    \item Multi-Task Learning(MTL多任务学习)
    
    \begin{enumerate}
        \item 多任务学习训练一个学习同时可以执行多个相关任务
        \item 在任务之间共享知识，提升各种任务性能。
    \end{enumerate}
    \item Multi-Modal Learning(MML多模态学习)
    \begin{enumerate}
        \item 整合来自多个模态的数据，通过结合互补信息源提升模型的理解能力。
    \end{enumerate}

\end{itemize}

\subsection{NLP的evaluation}
\begin{itemize}
    \item IR,TC,NER,POS等任务
    \begin{enumerate}
        \item Precision:所有检索到的item中，正确识别相关item所占的比例
        \item Recall:在数据中所有相关的item中，正确识别出相关item的比例
        \item F1 Score:精确率和召回率的Harmonic mean(调和平均数)
        \begin{itemize}
            \item Micro F1:对所有单个实例的精确值和召回率取平均值，不考虑类别大小，平等对待所有类别。
            \item Macro F1:分别计算每一个类别的F1平均分数并取平均值，平等对待所有类别。
        \end{itemize}
        \item Accuracy:正确预测数与总预测数的比例，适用于平衡的数据集
    \end{enumerate}

    \item MT,TS,GT等
    \begin{enumerate}
        \item BLEU(Bilingual Evaluation Understudy双语评估替补):衡量生成的翻译和参考翻译中n元语法的重叠度。
        \begin{itemize}
            \item 计算：unigram。BLEU2要取1-gram和2-gram的几何平均。
        \end{itemize}
        \item ROUGH(Recall-Oriented Understudy for Gisting Evaluation面向召回的要点评估替补指标):评估生成的摘要与参考摘要之间的n元语法、单词序列或句子的重叠程度。
        \begin{itemize}
            \item 如果匹配的单词有6个，在1-gram中有4个出现在人工参考文本之间，那么Rough1为$\frac{2}{3}$
        \end{itemize}
    \end{enumerate}
\end{itemize}

\section{Word Representation词汇表征}

\end{document}
